% Generated draft: requires mathematical and source review.
\documentclass[11pt]{article}
\usepackage[margin=25mm]{geometry}
\usepackage{fontspec}
\setmainfont{Latin Modern Roman}
\usepackage{amsmath,amssymb,mathtools,unicode-math}
\setmathfont{Latin Modern Math}
\usepackage{enumitem,needspace,xurl,hyperref}
\setlength{\parindent}{0pt}
\setlength{\parskip}{5pt}
\setlength{\emergencystretch}{3em}
\newcommand{\N}{\mathbb{N}}
\newcommand{\Z}{\mathbb{Z}}
\newcommand{\Q}{\mathbb{Q}}
\newcommand{\R}{\mathbb{R}}
\newcommand{\C}{\mathbb{C}}
\newcommand{\F}{\mathbb{F}}
\newcommand{\A}{\mathbb{A}}
\newcommand{\PP}{\mathbb{P}}
\newcommand{\E}{\mathbb{E}}
\newcommand{\eps}{\varepsilon}
\newcommand{\dd}{\,\mathrm{d}}
\begin{document}
\section*{UnsolvedMath ID: 20003181}
\textbf{Source problem code:} catalogue entry
\subsection{Definitions and mathematical statement}
Let \(S\) be a closed oriented surface of genus \(g\geq2\), \(\Gamma=\pi_1(S)\), and \(G=\operatorname{Sp}(4,\R)\). In a Toledo normalization with positive maximal value \(2g-2\), call \(\rho:\Gamma\to G\) positive maximal when \(\tau(\rho)=2g-2\). A normalization differing by a fixed factor must be adjusted consistently. The source does not explicitly give these surface and Toledo conventions; the negative maximal sign also requires its corresponding component labels.

Bill Goldman's question asks whether every maximal representation can be deformed to an endpoint whose image lies in a proper real Zariski-closed subgroup of \(G\). In a reductive character-variety formulation, this means a path of representation classes from \([\rho]\) to \([\rho_1]\), with \(\rho_1(\Gamma)\subset H\subsetneq G\) for some such subgroup \(H\). The subgroup may depend on the endpoint; the entire deformation is not required to lie in it. The exact representation-versus-character-variety convention requires verification.

The aid attributes to Bradlow, Garc\'ia-Prada, and Gothen a negative answer with an exact component characterization. In the indicated positive maximal convention, the exceptional components are
\[
 \mathcal R_c^0(S),\qquad c\in\Z,\quad0<c<2g-2,
\]
numbering \(2g-3\). Every representation in them is Zariski dense, and none can be deformed to a proper Zariski-closed subgroup. Every representation in every other maximal component admits the requested deformation. The integer invariant \(c\) and superscript \(0\) are labels from the cited classification; their intrinsic definitions are absent from the supplied record and are not invented here.

For every connected unramified cover \(p:S'\to S\) of degree \(m\geq1\), set \(g'=m(g-1)+1\). The aid reports
\[
 p^*(\mathcal R_c^0(S))\subseteq\mathcal R_{mc}^0(S'),\qquad
 \tau(p^*\rho)=m\tau(\rho).
\]
Since \(0<mc<2g'-2\), the pulled-back representation remains maximal and exceptional, so its non-deformability persists on every such finite cover. This is stronger than saying only that the original representation stays dense after restriction.

The aid also states that every finite-index restriction of an exceptional representation is Zariski dense and strongly irreducible in the standard real four-dimensional representation. Here strong irreducibility means irreducibility after every finite-index restriction. Its centralizer \emph{inside} \(\operatorname{Sp}(4,\R)\) is \(\{\pm I\}\); no identical centralizer assertion is made in the larger general linear group.

The original 2007 comments must remain historical context. Goldman says every maximal representation of ``the unitary group'' deforms as requested, without identifying that group. Anna Wienhard says the \(\operatorname{Sp}(2n,\R)\) problem should be easier for \(n>2\) but was then unknown, and records \(2(g-2)\) exotic components for \(n=2\) on which nothing was then known. This historical count is not silently identified with the aid's \(2g-3\) exceptional-component count; their terminology and scope need comparison. The extracted heading ``2 Geometry of the Hitchin component'' starts the next section and is not part of the question. The aid's later solved-in-literature assessment applies to its specified \(\operatorname{Sp}(4,\R)\) formulation, not automatically to every higher-rank or unitary variant.
\subsection{Short English statement}
The supplied literature answer is negative precisely in the stated exceptional maximal components, whose representations are all Zariski dense. The aid further reports that finite-cover pullback preserves the exceptional obstruction. Historical workshop counts and comments are retained separately, with component definitions and normalization still requiring source review.
\subsection{Sources}
\begin{itemize}
\item Local UnsolvedMath registry, permanent ID \(20003181\), exact source code \texttt{catalogue entry}; canonical \texttt{aim-topology-notes.json}, \texttt{notes[268]}, source tag \texttt{question}.
\item AIM workshop ``Representations of surface groups,'' 19--23 March 2007, Question 1.1 by Bill Goldman and Comments 1.2--1.3 by Bill Goldman and Anna Wienhard: \url{https://aimath.org/WWN/surfacegroups/surfacegroups.pdf}. The supplied verification calls the canonical record exact. No PDF was fetched.
\item The aid attributes the component theorem to Bradlow, Garc\'ia-Prada, and Gothen but supplies no paper title, identifier, or full authoritative reference; the original literature field is empty. The finite-cover corollary likewise has no complete citation in this record.
\item Research attempt 1: result \texttt{full\_solution}, problem status at run \texttt{solved}, rigor label \texttt{proved}, novelty confidence \texttt{low}; classification \texttt{SOLVED-IN-LITERATURE}, registry status \texttt{solved}. Created and updated 14 August 2026. These are supplied later assessments, not independent verification or a replacement of the 2007 historical comments.
\item Attribution: Ulam AI / UnsolvedMath Contributors, CC BY 4.0. The machine-generated research aid is not peer reviewed and requires independent verification before citation.
\end{itemize}
\end{document}
